\section{Beurteilung des Verlaufs des Aktienkurses}
\label{sec:aktienkurs}

Grundlage dieser Beurteilung sind die XETRA-Kurse aus
Abbildung~\ref{fig:kursverlauf} und Tabelle~\ref{tab:kursextrema}, deren
Herkunft und Pr\"ufung Abschnitt~\ref{subsec:kursverlauf} nennt, sowie die
Ergebnisse der Abschnitte~\ref{sec:jahresabschluss} und \ref{sec:strategie}.

% FAKTENBASIS 4
%  - Jahresschlusskurse: \KursSilvesterZFII, \KursSilvesterZFIII,
%    \KursSilvesterZFIV, \KursSilvesterZFV EUR (Abbildung \ref{fig:kursverlauf})
%  - Kursveraenderungen: \KursDeltaZFIII %, \KursDeltaZFIV %, \KursDeltaZFV %
%  - Jahreshoch/-tief (Tabelle \ref{tab:kursextrema}): 2024 Hoch \KursHochZFIV
%    EUR (Januar), Tief \KursTiefZFIV EUR (November) - Rueckgang \KursSpanneZFIV %
%  - 2025 Hoch \KursHochZFV EUR (Mai), Tief \KursTiefZFV EUR (Maerz)
%  - Auftragseingang 2023 \AuftragseingangZFIII, 2024 \AuftragseingangZFIV
%    Mio. EUR, Rueckgang \AuftragseingangDeltaZFIV %
%  - Umsatz 2024 \UmsatzDeltaZFIV %, operatives EBITDA \EbitdaOpDeltaZFIV %,
%    Konzernjahresueberschuss \NettogewinnDeltaZFIV % - alles Rekordwerte
%  - EBITDA-Marge 2024 \EbitdaMargeZFIV %, Dividendenrendite
%    \DividendenrenditeZFIV %
%  - Umsatzrueckgang 2025 \UmsatzDeltaZFV %, Marge dennoch ueber der Prognose
%    (Abschnitt \ref{subsec:ausblick}, Tabelle \ref{tab:prognose-ist})
%  - Klumpenrisiko E-Mobility \AnteilEMobility % (Abschnitt \ref{sec:strategie})
%  - Emissionspreis 2017 \IPOKurs EUR - ueber allen fuenf Schlusskursen
%  - Keine Analystenratings oder Kursziele belegbar (Abschnitt
%    \ref{subsec:analystenbotschaften})
\bk{Der Kursverlauf l\"asst sich aus den vorstehenden Analysen weitgehend
erkl\"aren, allerdings nicht anhand der Gewinn- und Verlustrechnung, sondern
anhand des Auftragseingangs. Er ist f\"ur einen Anlagenbauer der
Fr\"uhindikator der k\"unftigen Auslastung, w\"ahrend Umsatz und Ergebnis
abgearbeitete Auftr\"age abbilden. Dem Kursanstieg um \Pct{\KursDeltaZFIII}
im Jahr 2023 steht der Rekordauftragseingang von
\MioEUR{\AuftragseingangZFIII} gegen\"uber. Umgekehrt verlor die Aktie 2024
\PctBetrag{\KursDeltaZFIV}, obwohl Umsatz (\Pct{\UmsatzDeltaZFIV}),
operatives EBITDA (\Pct{\EbitdaOpDeltaZFIV}) und
Konzernjahres\-\"uberschuss (\Pct{\NettogewinnDeltaZFIV}) Rekordwerte
erreichten; erkl\"arungskr\"aftig ist allein der Auftragseingang, der um
\PctBetrag{\AuftragseingangDeltaZFIV} einbrach. Wie fr\"uh der Markt
reagierte, zeigt der unterj\"ahrige Verlauf: Das Jahreshoch von
\EURje{\KursHochZFIV} lag bereits im Januar 2024, also vor Ver\"offentlichung
der Rekordzahlen, das Tief von \EURje{\KursTiefZFIV} im November, ein Verlust
von \Pct{\KursSpanneZFIV} (Tabelle~\ref{tab:kursextrema}). Der Markt las die
Rekordzahlen offenbar als Abarbeitung eines auslaufenden Auftragsbestands,
was der Umsatzr\"uckgang von \PctBetrag{\UmsatzDeltaZFV} im Jahr 2025
best\"atigt hat. Dass der Kurs 2025 dennoch um \Pct{\KursDeltaZFV} zulegte,
passt dazu: Der R\"uckgang war seit dem Ausblick eingepreist, w\"ahrend die
Marge die eigene Prognose \"ubertraf (Tabelle~\ref{tab:prognose-ist}) und die
Nettoliquidit\"at trotz Dividende und R\"uckkauf zunahm. Nicht aus den
Ergebnissen erkl\"arbar bleibt das absolute Kursniveau: Alle f\"unf
Schlusskurse liegen deutlich unter dem Emissionspreis von \EURje{\IPOKurs}.
Darin spiegelt sich weniger die operative Entwicklung als die Bewertung der
Zyklusabh\"angigkeit, die beim E-Mobility-Anteil von rund
\Pct{\AnteilEMobility} unver\"andert hoch ist.}
